\begin{figure}[H]
\centering
\begin{tikzpicture}[>=Stealth,scale=0.92]
  \draw[rounded corners=3pt,fill=gray!8,draw=gray!60]
    (-3.3,-1.55) rectangle (2.65,1.55);
  \node[above] at (-0.3,1.6) {telescopio};

  \draw[fill=blue!8,draw=blue!65!black,thick] (-5.8,-1.65) rectangle (-5.55,1.65);
  \node[align=center,left] at (-5.85,0) {fuente extensa\\uniforme};
  \foreach \y in {-1.0,-0.5,0,0.5,1.0}
    \draw[orange!85!black,thick,->] (-5.5,\y) -- (-3.45,\y);

  \draw[<->,very thick,blue!75!black] (-3.05,-1.35) -- (-3.05,1.35);
  \node[below] at (-3.05,-1.55) {objetivo};
  \draw[<->,very thick,cyan!65!blue] (2.25,-0.82) -- (2.25,0.82);
  \node[below] at (2.25,-1.55) {ocular};
  \draw[dashed,gray!75] (-3.45,0) -- (6.1,0);

  \draw[blue!70!black,thick] (-3.05,1.15) -- (2.25,0.48) -- (4.35,0.16);
  \draw[blue!70!black,thick] (-3.05,-1.15) -- (2.25,-0.48) -- (4.35,-0.16);
  \draw[blue!55,thick] (-3.05,0.45) -- (2.25,0.18) -- (4.35,0.06);
  \draw[blue!55,thick] (-3.05,-0.45) -- (2.25,-0.18) -- (4.35,-0.06);

  \draw[fill=white,draw=gray!70,thick] (4.35,-1.45) rectangle (4.47,1.45);
  \fill[yellow!75!orange,draw=orange!80!black] (4.29,-0.16) rectangle (4.53,0.16);
  \draw[<->] (4.85,-0.16) -- (4.85,0.16)
    node[midway,right,fill=white,inner sep=1.5pt] {$d_{\mathrm{exp}}$};
  \node[above,align=center] at (4.41,1.55) {pantalla móvil\\en el plano de la pupila};
  \draw[<->,gray!80] (3.55,-1.85) -- (5.25,-1.85);
  \node[below] at (4.4,-1.86) {ajuste longitudinal};

  \draw[thick] (5.35,-0.55) -- (6.15,-0.55);
  \foreach \x in {5.35,5.45,...,6.15}
    \draw (\x,-0.55) -- (\x,-0.42);
  \node[below] at (5.75,-0.62) {escala};
\end{tikzpicture}
\caption{Montaje para localizar y medir experimentalmente la pupila de salida. La pantalla se desplaza hasta obtener el disco luminoso más nítido.}
\end{figure}